\ea\label{ex:text11-1}
Amaa humhum na inu gorogoroho vaa Teapu\\
\gll a=maa	hum{$\sim$}hum	n=a	inu	goro{$\sim$}goroho	vaa	teapu\\
     \textsc{art}2.\textsc{sg}=\textsc{plm}	\textsc{red}{$\sim$}part	3\textsc{sg}.\textsc{poss}=\textsc{art}2.\textsc{sg}	house	\textsc{red}{$\sim$}sleep	\textsc{lnk}	Teop.reagion\\
\glt `Some parts of the Teop sleeping house'
\z

\ea\label{ex:text11-2}
TOON\\
\gll o	toon\\
     \textsc{art}3.\textsc{sg}	post\\
\glt `POST'
\z

\ea\label{ex:text11-3}
A ta naono vai a ta beera, bebeahu ae kikis to toon riori bona inu,\\
\gll a	ta	naono	vai	a	ta	beera	bebeahu	ae	kikis	to	toon	ri=ori	bona	inu\\
     \textsc{art}2.\textsc{sg}	piece	wood	\textsc{dem}6	\textsc{art}2.\textsc{sg}	piece	big	long	and1	strong	\textsc{rel}	use.as.post.of	3\textsc{pl}.\textsc{ipfv}=3\textsc{pl}.\textsc{pron}	\textsc{obj}.\textsc{art}2.\textsc{sg}	house\\
\glt `A big, long and strong pole that is used as a post of the house,'
\z

\ea\label{ex:text11-4}
ge eve o naono to dee maana bona inu.\\
\gll ge	eve	o	naono	to	dee	maa=na	bona	inu\\
     or	3\textsc{sg}.\textsc{pron}	\textsc{art}3.\textsc{sg}	pole	\textsc{rel}	carry	hither=3\textsc{pl}.\textsc{ipfv}	\textsc{obj}.\textsc{art}2.\textsc{sg}	house\\
\glt `or, (put differently,) it is a pole that carries the house.'
\z

\ea\label{ex:text11-5}
Bea ta tovihi, na tei koana tea beera na inu.\\
\gll be=a	ta	tovihi	na	tei	koa=na	te=a	beera	n=a	inu\\
     \textsc{compl}2=\textsc{art}2.\textsc{sg}	piece	how.many	\textsc{tm}2	be.there	just=3\textsc{sg}.\textsc{ipfv}	\textsc{prep}=\textsc{art}2.\textsc{sg}	size	3\textsc{sg}.\textsc{poss}=\textsc{art}2.\textsc{sg}	house\\
\glt `How many, (that) depends on the size of the house.'
\z

\ea\label{ex:text11-6}
Evehee sa gonogono vahogo haari,\\
\gll evehee	sa	gono{$\sim$}gono	vahogo	haa=ri\\
     but	\textsc{neg}1	\textsc{red}{$\sim$}get	carelessly	\textsc{neg}2=3\textsc{pl}.\textsc{ipfv}\\
\glt `But they do not get them in a careless way,'
\z

\newpage
\ea\label{ex:text11-7}
ahiki, o naono to mataa ki nana bono toon o toraha.\\
\gll ahiki	o	naono	to	mataa	ki	nana	bono	toon	o	toraha\\
     no	\textsc{art}3.\textsc{sg}	tree	\textsc{rel}	good	\textsc{dat}	3\textsc{sg}.\textsc{ipfv}	\textsc{obj}.\textsc{art}3.\textsc{sg}	post	\textsc{art}3.\textsc{sg}	New.Guinea.Teak\\
\glt `no, the tree that is good for a post is the New Guinea Teak.'
\z

\ea\label{ex:text11-8}
O naono o kikis ae o vaanagu koava.\\
\gll o	naono	o	kikis	ae	o	vaanagu	koava\\
     \textsc{art}3.\textsc{sg}	wood	\textsc{art}3.\textsc{sg}	strong	and1	\textsc{art}3.\textsc{sg}	resist	rot\\
\glt `The wood is strong and resistant to rotting.'
\z

\ea\label{ex:text11-9}
PAEE\\
\gll o	paee\\
     \textsc{art}3.\textsc{sg}	horizontal.beam\\
\glt `HORIZONTAL BEAM (joist, tie beam, joist, and the main ridge pole)'
\z

\ea\label{ex:text11-10}
O naono o bebeahu to vaobete riori toon na maa paos.\\
\gll o	naono	o	bebeahu	to	va-obete	ri=ori	toon	n=a	maa	paos\\
     \textsc{art}3.\textsc{sg}	beam	\textsc{art}3.\textsc{sg}	long	\textsc{rel}	\textsc{caus}-lie	3\textsc{pl}.\textsc{ipfv}=3\textsc{pl}.\textsc{pron}	top	3\textsc{sg}.\textsc{poss}=\textsc{art}3.\textsc{sg}	\textsc{plm}	post\\
\glt `The long beam that is put on top of the posts.'
\z

\ea\label{ex:text11-11}
Evehee toro naono vai haabana o kikis, tabae eve to kahi dee nana bono paranda.\\
\gll evehee	toro	naono	vai	haabana	o	kikis	tabae	eve	to	kahi	dee	nana	bono	paranda\\
     but	\textsc{tm}7	beam	\textsc{dem}6	again	\textsc{art}3.\textsc{sg}	strong	because	3\textsc{sg}.\textsc{pron}	\textsc{rel}	\textsc{tm}1	carry	3\textsc{sg}.\textsc{ipfv}	\textsc{obj}.\textsc{art}3.\textsc{sg}	floor\\
\glt `But (it) must again be a strong beam because it will carry the floor.'\footnote{lit. `... because it (is the one) that will carry the floor.'}
\z

\ea\label{ex:text11-12}
A ta tovihi? Na tei koana tea bero tea toon.\\
\gll a	ta	tovihi	na	tei	koa=na	te=a	bero	te=a	toon\\
     \textsc{art}2.\textsc{sg}	piece	how.many	\textsc{tm}2	be.there	just=3\textsc{sg}.\textsc{ipfv}	\textsc{prep}=\textsc{art}2.\textsc{sg}	quantity	\textsc{prep}=\textsc{art}3.\textsc{pl}	post\\
\glt `How many? That just depends on the number of the posts.'
\z

\ea\label{ex:text11-13}
O hum to teitei bata nae mee ei toon na maa sunsun.\\
\gll o	hum	to	tei{$\sim$}tei	bata	na=e	mee	ei	toon	n=a	maa	sunsun\\
     \textsc{art}3.\textsc{sg}	place	\textsc{rel}	\textsc{red}{$\sim$}stay	\textsc{cont}	3\textsc{sg}.\textsc{ipfv}=3\textsc{sg}.\textsc{ipfv}	also	\textsc{dem}2	top	3\textsc{sg}.\textsc{poss}=\textsc{art}2.\textsc{sg}	\textsc{plm}	standing.post\\
\glt `A place where they also lie, this is on top of the (posts called) sunsun, standing posts.'
\z

\ea\label{ex:text11-14}
SUNSUN\\
\gll sunsun\\
     standing.post\\
\glt `STANDING POST'
\z

\ea\label{ex:text11-15}
O sunsun o naono to sun nana toon na maa paee,\\
\gll o	sunsun	o	naono	to	sun	nana	toon	n=a	maa	paee\\
     \textsc{art}3.\textsc{sg}	standing.post	\textsc{art}3.\textsc{sg}	pole	\textsc{rel}	stand	3\textsc{sg}.\textsc{ipfv}	top	3\textsc{sg}.\textsc{poss}=\textsc{art}2.\textsc{sg}	\textsc{plm}	horizontal.beam\\
\glt `The standing post is a pole that stands on the horizontal beams.'
\z

\ea\label{ex:text11-16}
O meho paee repaa obete taono vaha ni maa bona,\\
\gll o	meho	paee	re=paa	obete	taono	vaha	ni	maa	bona\\
     \textsc{art}3.\textsc{sg}	other3.\textsc{sg}	horizontal.beam	\textsc{consec}=\textsc{tm}3	lie	put.on.top	again	\textsc{appl}	hither	4\textsc{sg}.\textsc{pron}\\
\glt `And another horizontal beam lies on it,'
\z

\ea\label{ex:text11-17}
O varuvaru repaa obete maa toon nae.\\
\gll o	varuvaru	re=paa	obete	maa	toon	na=e\\
     \textsc{art}3.\textsc{sg}	rafter	\textsc{consec}=\textsc{tm}3	lie	hither	top	3\textsc{sg}.\textsc{poss}=3\textsc{sg}.\textsc{pron}\\
\glt `And then the rafter lies on top of it.'
\z

\ea\label{ex:text11-18}
Eve o naono to kahi ato ni nana bono hum inu vaa gina\\
\gll eve	o	naono	to	kahi	ato	ni	nana	bono	hum	inu	vaa	gina\\
     3\textsc{sg}.\textsc{pron}	\textsc{art}3.\textsc{sg}	pole	\textsc{rel}	\textsc{tm}1	hold	\textsc{appl}	3\textsc{sg}.\textsc{ipfv}	\textsc{obj}.\textsc{art}3.\textsc{sg}	part	house	\textsc{lnk}	top\\
\glt `It is the pole that holds the upper part of the house.'
\z

\ea\label{ex:text11-19}
GAASUU\\
\gll gaasuu\\
     kingpost\\
\glt `KING POST'
\z

\ea\label{ex:text11-20}
O gaasuu o naono to vavaasun riori tobina inu toon no paee.\\
\gll o	gaasuu	o	naono	to	va{$\sim$}vaa-sun	ri=ori	tobina	inu	toon	n=o	paee\\
     \textsc{art}3.\textsc{sg}	kingpost	\textsc{art}3.\textsc{sg}	pole	\textsc{rel}	\textsc{red}{$\sim$}\textsc{caus}-stand	3\textsc{pl}.\textsc{ipfv}=3\textsc{pl}.\textsc{pron}	middle	house	top	3\textsc{sg}.\textsc{poss}=\textsc{art}3.\textsc{sg}	horizontal.beam\\
\glt `The gaasu is the pole that they stand in the middle of the house on top of the tie beam.'
\z

\ea\label{ex:text11-21}
Beo tovihi? Na tei koana tea bebeahu na inu.\\
\gll be=o	tovihi	na	tei	koa=na	te=a	bebeahu	n=a	inu\\
     when=\textsc{art}3.\textsc{sg}	how.many	\textsc{tm}2	be.there	just=3\textsc{sg}.\textsc{ipfv}	\textsc{prep}=\textsc{art}2.\textsc{sg}	length	3\textsc{sg}.\textsc{poss}=\textsc{art}2.\textsc{sg}	house\\
\glt `How many? That just depends on the length of the house.'
\z

\ea\label{ex:text11-22}
Evehee a kiu tea maa naono bari mene dee\\
\gll evehee	a	kiu	te=a	maa	naono	bari	mene	dee\\
     but	\textsc{art}2.\textsc{sg}	work	\textsc{prep}=\textsc{art}2.\textsc{sg}	\textsc{plm}	beam	\textsc{ana}.\textsc{pl}	for	carry\\
\glt `The function of these poles is to carry (the main ridge pole)'
\z

\ea\label{ex:text11-23}
ge eve to vaatoka maa riori bono peho paee.\\
\gll ge	eve	to	vaa-toka	maa	ri=ori	bono	peho	paee\\
     or	3\textsc{sg}.\textsc{pron}	\textsc{rel}	\textsc{caus}-perch	hither	3\textsc{pl}.\textsc{ipfv}=3\textsc{pl}.\textsc{pron}	\textsc{obj}.\textsc{art}3.\textsc{sg}	one3.\textsc{sg}	horizontal.beam\\
\glt `or, (in other words,) it is the one on which they make the main ridge pole perch.'
\z

\ea\label{ex:text11-24}
Amaa ta naono bari mee amaa ta kakato koa.\\
\gll a=maa	ta	naono	bari	mee	a=maa	ta	kakato	koa\\
     \textsc{art}2.\textsc{sg}=\textsc{plm}	piece	wood	\textsc{ana}.\textsc{pl}	also	\textsc{art}2.\textsc{sg}=\textsc{plm}	piece	short	quite\\
\glt `These pieces of wood are quite short.'
\z

\newpage
\ea\label{ex:text11-25}
VARUVARU\\
\gll varuvaru\\
     rafter\\
\glt `RAFTER'
\z

\ea\label{ex:text11-26}
A varuvaru eve a naono to kahi vaobete mau ori toon no paee vai\\
\gll a	varuvaru	eve	a	naono	to	kahi	va-obete	ma=u	ori	toon	n=o	paee	vai\\
     \textsc{art}3.\textsc{pl}	rafter	3\textsc{sg}.\textsc{pron}	\textsc{art}3.\textsc{pl}	beam	\textsc{rel}	\textsc{tm}1	\textsc{caus}-lie	hither=\textsc{pfv}	3\textsc{pl}.\textsc{pron}	top	3\textsc{sg}.\textsc{poss}=\textsc{art}3.\textsc{sg}	horizontal.beam	\textsc{dem}6\\
\glt `The rafters are the beams that are laid on the horizontal beam'
\z

\ea\label{ex:text11-27}
to paa vaobete vori toon na mata gaasuu.\\
\gll to	paa	va-obete	v=ori	toon	n=a	ma=ta	gaasuu\\
     \textsc{rel}	\textsc{tm}3	\textsc{caus}-lie	\textsc{pfv}=3\textsc{pl}.\textsc{pron}	top	3\textsc{sg}.\textsc{poss}=\textsc{art}2.\textsc{sg}	\textsc{plm}=piece	kingpost\\
\glt `that is laid on to the kingposts.'
\z

\ea\label{ex:text11-28}
Eve mee a naono to kahi vaave taono riori bona atovo.\\
\gll eve	mee	a	naono	to	kahi	vaave	taono	ri=ori	bona	atovo\\
     3\textsc{sg}.\textsc{pron}	also	\textsc{art}3.\textsc{pl}	pole	\textsc{rel}	\textsc{tm}1	fix.thatch	put.on.top	3\textsc{pl}.\textsc{ipfv}=3\textsc{pl}.\textsc{pron}	\textsc{obj}.\textsc{art}3.\textsc{pl}	sago.palm.thatch\\
\glt `They are also the beams on which the sago palm thatches will be fixed.'
\z

\ea\label{ex:text11-29}
A bero nae mee na tei habana koana tea bebeahu na inu.\\
\gll a	bero	na=e	mee	na	tei	habana	koa=na	te=a	bebeahu	n=a	inu\\
     \textsc{art}2.\textsc{sg}	quantity	3\textsc{sg}.\textsc{poss}=3\textsc{sg}.\textsc{pron}	also	\textsc{tm}2	be.there	again	just=3\textsc{sg}.\textsc{ipfv}	\textsc{prep}=\textsc{art}2.\textsc{sg}	length	3\textsc{sg}.\textsc{poss}=\textsc{art}2.\textsc{sg}	house\\
\glt `Their number also depends only on the length of the house.'
\z

\ea\label{ex:text11-30}
A naono nae a vahavahara naono vaa tea apao.\\
\gll a	naono	na=e	a	vaha{$\sim$}vahara	naono	vaa	te=a	apao\\
     \textsc{art}3.\textsc{pl}	tree	3\textsc{sg}.\textsc{poss}=3\textsc{sg}.\textsc{pron}	\textsc{art}3.\textsc{pl}	\textsc{red}{$\sim$}little.\textsc{pl}	tree	\textsc{lnk}	\textsc{prep}=\textsc{art}2.\textsc{sg}	old.garden\\
\glt `The trees used for them are little trees from the old garden.'
\z

\ea\label{ex:text11-31}
A varuvaru he toro tapaku tea buaku a papana tea papa matau ae papa keruka\\
\gll a	varuvaru	he	toro	ta-paku	te=a	buaku	a	papana	te=a	papa	matau	ae	papa	keruka\\
     \textsc{art}3.\textsc{pl}	rafter	\textsc{advers}	\textsc{tm}7	\textsc{acaus}	\textsc{prep}=\textsc{art}2.\textsc{sg}	two	\textsc{art}2.\textsc{sg}	side	\textsc{prep}=\textsc{art}2.\textsc{sg}	side	right	and1	side	left\\
\glt `The rafters must be made on the two sides, on the right side and the left side'
\z

\ea\label{ex:text11-32}
ore paa vatasu vatatana tobina toon no paee vai\\
\gll o=re	paa	va-tasu	va-tatana	tobina	toon	n=o	paee	vai\\
     3.\textsc{pron}=\textsc{consec}	\textsc{tm}3	\textsc{mult}-hit	\textsc{mult}-meet	middle	top	3\textsc{sg}.\textsc{poss}=\textsc{art}3.\textsc{sg}	horizontal.beam	\textsc{dem}6\\
\glt `so that they touch each other in the middle (of the roof) on top of the horizontal beam'
\z

\ea\label{ex:text11-33}
to paa vaatoka vori toon na maa gaasuu vaa tobina.\\
\gll to	paa	vaa-toka	v=ori	toon	n=a	maa	gaasuu	vaa	tobina\\
     \textsc{rel}	\textsc{tm}3	\textsc{caus}-perch	\textsc{pfv}=3\textsc{pl}.\textsc{pron}	top	3\textsc{sg}.\textsc{poss}=\textsc{art}2.\textsc{sg}	\textsc{plm}	kingpost	\textsc{lnk}	middle\\
\glt `that they have placed on the kingposts in the middle.'
\z

\ea\label{ex:text11-34}
VAATETEKOIO\\
\gll vatetekoio\\
     upper.ridgepole\\
\glt `Upper ridgepole'
\z

\ea\label{ex:text11-35}
O naono bebeahu to kahi vaatoka vaha maa riori toon na maa varuvaru vai to paa vatasu vatatana vuru.\\
\gll o	naono	bebeahu	to	kahi	vaa-toka	vaha	maa	ri=ori	toon	n=a	maa	varuvaru	vai	to	paa	va-tasu	va-tatana	vuru\\
     \textsc{art}3.\textsc{sg}	beam	long	\textsc{rel}	\textsc{tm}1	\textsc{caus}-perch	again	hither	3\textsc{pl}.\textsc{ipfv}=3\textsc{pl}.\textsc{pron}	top	3\textsc{sg}.\textsc{poss}=\textsc{art}2.\textsc{sg}	\textsc{plm}	rafter	\textsc{dem}6	\textsc{rel}	\textsc{tm}3	\textsc{mult}-hit	\textsc{mult}-meet	earlier\\
\glt `A long beam that will be placed on top of the rafters that have touched each other (on the lower ridgepole).'
\z

\ea\label{ex:text11-36}
Eve to kahi ato vakikis habana nana tea inu.\\
\gll eve	to	kahi	ato	va-kikis	habana	nana	te=a	inu\\
     3\textsc{sg}.\textsc{pron}	\textsc{rel}	\textsc{tm}1	hold	\textsc{advr}-strong	again	3\textsc{sg}.\textsc{ipfv}	\textsc{prep}=\textsc{art}2.\textsc{sg}	house\\
\glt `It will also hold the house firmly together.'
\z

\ea\label{ex:text11-37}
Evehee a kiu teo naono bona rakaha, ei meno kovuaka.\\
\gll evehee	a	kiu	te=o	naono	bona	rakaha	ei	meno	kovuaka\\
     but	\textsc{art}2.\textsc{sg}	work	\textsc{prep}=\textsc{art}3.\textsc{sg}	beam	\textsc{ana}.\textsc{sg}	in.fact	\textsc{dem}2	for.\textsc{art}3.\textsc{sg}	top.thatch.sheet\\
\glt `But the main function of this beam, this is for the top thatch sheet.'
\z

\ea\label{ex:text11-38}
Eve to kahi vaobete maa riori bono kovuaka ore paa obete vamataa koa maa teebona.\\
\gll eve	to	kahi	va-obete	maa	ri=ori	bono	kovuaka	o=re	paa	obete	va-mataa	koa	maa	teebona\\
     3\textsc{sg}.\textsc{pron}	\textsc{rel}	\textsc{tm}1	\textsc{caus}-lie	hither	3\textsc{pl}.\textsc{ipfv}=3\textsc{pl}.\textsc{pron}	\textsc{obj}.\textsc{art}3.\textsc{sg}	top.thatch.sheet	3.\textsc{pron}=\textsc{consec}	\textsc{tm}3	lie	\textsc{advr}-good	just	hither	there\\
\glt `It is (the beam) where they will put the top thatch sheet so that it properly lies there.'
\z

\ea\label{ex:text11-39}
HOGO\\
\gll hogo\\
     floor.joist\\
\glt `FLOOR JOIST'
\z

\ea\label{ex:text11-40}
O naono to hogo vaha ni maa riori paana, be obete bata maana toon na maa paee vai tea maa toon.\\
\gll o	naono	to	hogo	vaha	ni	maa	ri=ori	paa-na	be	obete	bata	maa=na	toon	n=a	maa	paee	vai	te=a	maa	toon\\
     \textsc{art}3.\textsc{sg}	beam	\textsc{rel}	use.as.floor.joist	again	\textsc{appl}	hither	3\textsc{pl}.\textsc{ipfv}=3\textsc{pl}.\textsc{pron}	under-\textsc{derel}	while	lie	\textsc{cont}	hither=3\textsc{sg}.\textsc{ipfv}	top	3\textsc{sg}.\textsc{poss}=\textsc{art}2.\textsc{sg}	\textsc{plm}	horizontal.beam	\textsc{dem}6	\textsc{prep}=\textsc{art}2.\textsc{sg}	\textsc{plm}	post\\
\glt `The beams that they put underneath (the floor), while they lie on the horizontal beams (that lie) on the posts.'
\z

\ea\label{ex:text11-41}
Eve to kahi obete mau o paranda.\\
\gll eve	to	kahi	obete	ma=u	o	paranda\\
     3\textsc{sg}.\textsc{pron}	\textsc{rel}	\textsc{tm}1	lie	hither=\textsc{pfv}	\textsc{art}3.\textsc{sg}	floor\\
\glt `They are the ones on which the floor will lie.'
\z

\ea\label{ex:text11-42}
O tovihi? Na tei pete koana tea bebeahu ae tea tatabonata na inu.\\
\gll o	tovihi	na	tei	pete	koa=na	te=a	bebeahu	ae	te=a	tatabonata	n=a	inu\\
     \textsc{art}3.\textsc{sg}	how.many	\textsc{tm}2	be.there	also	just=3\textsc{sg}.\textsc{ipfv}	\textsc{prep}=\textsc{art}2.\textsc{sg}	length	and1	\textsc{prep}=\textsc{art}2.\textsc{sg}	width	3\textsc{sg}.\textsc{poss}=\textsc{art}2.\textsc{sg}	house\\
\glt `How many? That also just depends on the length and the width of the house.'
\z
